
opérations \(\psi^k\)
définissent d'après 7.3 des endomorphismes additifs du foncteur en \(\lambda\)-anneaux \(\widehat G\)~; d'après 1.10, ils sont caractérisés par leur valeur pour \(1+Xt\). On a :
\[
  \frac{d}{dt}(\lambda_u(1+Xt))/\lambda_u(1+Xt))
  =
  \frac{1+Xt}{\underline{1}+\{1+Xt\}u}
  =
  \sum_{k=1}^{\infty}(-1)^{k-1}\{1+Xt\}^{\circ k}u^{k-1}.
\]
Par suite :
\begin{equation}\tag{7.4.1}
  \psi^k(1+Xt)=1+X^k t.
\end{equation}

\textbf{Proposition 7.5.} i) \uline{Si} \(K\) \uline{est un} \(\lambda\)-\uline{anneau}, \uline{les opérations} \(\psi^k\) \uline{vérifient} \uline{les identités} :
\begin{equation}\tag{7.5.1}
  \psi^k(1)=1~;\qquad
  \psi^k(xy)=\psi^k(x)\psi^k(y)~;\qquad
  \psi^k\circ\psi^{k'}=\psi^{kk'}.
\end{equation}
ii) \uline{Réciproquement}, \uline{si les} \(\psi^t\) \uline{vérifient ces identités}, \uline{et si} \(K\) \uline{est de} \uline{caractéristique nulle}, \(K\) \uline{est un} \(\lambda\)-\uline{anneau}.

i) Si \(K\) est un \(\lambda\)-anneau, on a \(\lambda_t(1)=1+t\), d'où : \(\frac{d}{dt}(\lambda_t(1))/\lambda_t(1)=1/1+t=\sum_{k=1}^{\infty}(-1)^{k-1}t^{k-1}\), et par suite :
\[
  \psi^k(1)=1.
\]

Pour montrer la seconde des relations (7.5.1), on considère le diagramme :
\begin{equation}\tag{7.5.2}
\begin{tikzcd}[column sep=large,row sep=large]
K\times K \arrow[r,"\lambda_t\times\lambda_t"] \arrow[d] &
\widehat G(K)\times\widehat G(K) \arrow[r,"\Phi^k\times\Phi^k"] \arrow[d,"\circ"'] &
K\times K \arrow[d]\\
K \arrow[r,"\lambda_t"] &
\widehat G(K) \arrow[r,"\Phi^k"] &
K ,
\end{tikzcd}
\end{equation}
où les flèches verticales sont définies par la multiplication, et les applications \(\Phi^k\) comme suit : l'application « dérivée logarithmique » est
